\subsection{Galois Algebras}

Let K be a commutative field. If E is a K-algebra, then we denote its automorphism group by \(\operatorname{Aut}_{\mathrm{K}}(\mathrm{E})\) . If E is a Galois extension of the field K, then the group \(\operatorname{Aut}_{\mathrm{K}}(\mathrm{E})\) is simply the Galois group \(\operatorname{Gal}(\mathrm{E}/\mathrm{K})\) (V, §10, No.~2, p.~58).

Let G be a group. A \((K,G)\) -algebra is a K-algebra E endowed with a group homomorphism \(\lambda:G\to\operatorname{Aut}_{K}(E)\) . The homomorphism \(\lambda\) then endows E with the structure of a group with operators in G as well as the structure of a left K{[}G{]}-module with external law given by

\[
\Bigl (\sum_ {g \in \mathrm{G}} \mu_ {g} g \Bigr) x = \sum_ {g \in \mathrm{G}} \mu_ {g} \lambda (g). x\tag{14}
\]

for every \(x \in \mathrm{L}\) and every element \((\mu_g)_{g \in \mathrm{G}}\) of \(\mathrm{K}[\mathrm{G}]\) . A morphism of \((\mathrm{K}, \mathrm{G})\) -algebras is a morphism of algebras that is also a morphism of groups with operators.

For every family \((\mathrm{E}_{i})_{i\in\mathrm{I}}\) of \((\mathrm{K},\mathrm{G})\) -algebras, the product K-algebra \(\prod_{i\in I}E_{i}\) endowed with its structure of a product group with operators is a \((\mathrm{K},\mathrm{G})\) -algebra.

If \(\mathrm{E}\) is a (K, G)-algebra, then the set \(\mathrm{E}^{\mathrm{G}}\) of elements of \(\mathrm{E}\) invariant under \(\mathrm{G}\) is a subalgebra of \(\mathrm{E}\) .

Let E be a (K, G)-algebra, where G acts by \(\lambda: G \to \operatorname{Aut}_{\mathrm{K}}(\mathrm{E})\) . If \(K'\) is an extension of K, then for any \(g \in G\) , let \(\lambda'(g)\) be the automorphism \(\operatorname{Id}_{\mathrm{K}'} \otimes \lambda(g)\) of the \(K'\) -algebra \(\operatorname{L}_{(\mathrm{K}')}\) . Then \(\lambda': g \mapsto \lambda'(g)\) endows \(\operatorname{L}_{(\mathrm{K}')}\) with the structure of a \((\mathrm{K}', \mathrm{G})\) -algebra.

Given a group H and H-sets X and Y, we denote by \(\mathcal{F}_{\mathrm{H}}(\mathrm{X},\mathrm{Y})\) the set of morphisms of H-sets from X to Y. It is the set of mappings \(f:X\to Y\) such that \(f(hx)=hf(x)\) for every \(h\in H\) and \(x\in X\) .

Let G be a group with identity element e, let H be a subgroup of G, and let E be a (K,H)-algebra. The (K,G)-algebra deduced from E by coinduction from H to G, denoted by \(\mathrm{Coind}_{\mathrm{H}}^{\mathrm{G}}(\mathrm{E})\) , is the K-algebra \(\mathcal{F}_{\mathrm{H}}(\mathrm{G},\mathrm{E})\) endowed with the action of G given by the homomorphism \(\lambda\) from G to \(\mathrm{Aut}_{\mathrm{K}}(\mathrm{Coind}_{\mathrm{H}}^{\mathrm{G}}(\mathrm{E}))\) defined by

\[
(\lambda (g) \cdot f) (g ^ {\prime}) = f (g ^ {\prime} g)\tag{15}
\]

for \(f \in \text{Coind}_{\text{H}}^{\text{G}}(\text{E})\) and \(g, g' \in G\) .

Lemma 7. --- a) Let S be a subset of G such that every element of G can be written uniquely as hs with \(h \in H\) and \(s \in S\) . Then the mapping that sends f to \((f(s))_{s \in S}\) is an isomorphism from the K-algebra \(\operatorname{Coind}_{\mathrm{H}}^{\mathrm{G}}(\mathrm{E})\) to the K-algebra \(\mathcal{F}(\mathrm{S}, \mathrm{E})\) of mappings from S to E.

\begin{enumerate}
\def\labelenumi{\alph{enumi})}
\setcounter{enumi}{1}
\tightlist
\item
  The algebra \(\operatorname{Coind}_{\mathrm{H}}^{\mathrm{G}}(\mathrm{E})\) has finite degree over K if and only if E has finite degree over K and the index of H in G is finite. In this case, we have the formula
\end{enumerate}

\[
[ \mathrm{Coind} _ {\mathrm{H}} ^ {\mathrm{G}} (\mathrm{E}): \mathrm{K} ] = (\mathrm{G}: \mathrm{H}) [ \mathrm{E}: \mathrm{K} ].
\]

\begin{enumerate}
\def\labelenumi{\alph{enumi})}
\setcounter{enumi}{2}
\tightlist
\item
  Let \(E^{H}\) be the algebra of invariants of the group H in E and \(\operatorname{Coind}_{\mathrm{H}}^{\mathrm{G}}(\mathrm{E})^{\mathrm{G}}\) that of the group G in \(\operatorname{Coind}_{\mathrm{H}}^{\mathrm{G}}(\mathrm{E})\) . The mapping \(f \mapsto f(e)\)
\end{enumerate}

from \(\mathrm{Coind}_{\mathrm{H}}^{\mathrm{G}}(\mathrm{E})\) to E restricts to an algebra isomorphism from \(\mathrm{Coind}_{\mathrm{H}}^{\mathrm{G}}(\mathrm{E})^{\mathrm{G}}\) to \(\mathrm{E}^{\mathrm{H}}\) .

Assertion a) follows from the definitions and implies b). By (15), a mapping from G to E is an element of \(\mathrm{Coind}_{\mathrm{H}}^{\mathrm{G}}(\mathrm{E})\) invariant under G if and only if it is constant with value an element of E invariant under H.

Remark 1. --- Let G be a group, H a subgroup of G, and N a subgroup of H. Let E be a (K,N)-algebra. Let \(\alpha\) be an element of \(\mathcal{F}_{\mathrm{H}}(\mathrm{G},\mathcal{F}_{\mathrm{N}}(\mathrm{H},\mathrm{E}))\) . We have the relation

\[
\alpha (g) (n h) = n (\alpha (g) (h))\tag{16}
\]

for every \(g \in \mathbf{G}\) , \(h \in \mathbf{H}\) , and \(n \in \mathbf{N}\) and the relations

\[
\alpha (h g) (h ^ {\prime}) = (h \alpha (g)) (h ^ {\prime}) = \alpha (g) (h ^ {\prime} h)\tag{17}
\]

for all \(h, h' \in H\) and every \(g \in G\) . Consequently, \(\alpha(ng)(e) = \alpha(g)(n) = n(\alpha(g)(e))\) for every \(g \in G\) and \(n \in N\) . We can therefore consider the mapping

\[
\psi : \mathcal {F} _ {\mathrm{H}} (\mathrm{G}, \mathcal {F} _ {\mathrm{N}} (\mathrm{H}, \mathrm{E})) \to \mathcal {F} _ {\mathrm{N}} (\mathrm{G}, \mathrm{E})
\]

defined by the relation \(\psi(\alpha)(g)=\alpha(g)(e)\) for \(\alpha\) in \(\mathcal{F}_{\mathrm{H}}(\mathrm{G},\mathcal{F}_{\mathrm{N}}(\mathrm{H},\mathrm{E}))\) and g in G. The mapping \(\psi\) is an algebra isomorphism from \(\operatorname{Coind}_{\mathrm{H}}^{\mathrm{G}}(\operatorname{Coind}_{\mathrm{N}}^{\mathrm{H}}(\operatorname{E}))\) to \(\operatorname{Coind}_{\mathrm{N}}^{\mathrm{G}}(\operatorname{E})\) whose inverse sends an element \(\beta\) from \(\mathcal{F}_{\mathrm{N}}(\mathrm{G},\mathrm{E})\) to the mapping \(\alpha\) defined by the relation \(\alpha(g)(h)=\beta(hg)\) for \(g\in G\) and \(h\in H\) .

We now suppose given a finite group G and a reduced (V, §6, No.~6, p.~32) commutative K-algebra L of finite degree endowed with an action of G given by a homomorphism \(\lambda\) from G to \(\mathrm{Aut}_{\mathrm{K}}(\mathrm{L})\) . For \(x \in \mathrm{L}\) and \(g \in \mathrm{G}\) , we denote by \(g \cdot x\) the transform of \(x\) under the automorphism \(\lambda(g)\) of L. Let \(\mathcal{S}\) be the set of maximal ideals of L; we denote by \(g \cdot \mathfrak{m}\) the transform of an element \(\mathfrak{m}\) of \(\mathcal{S}\) under the automorphism \(\lambda(g)\) of L. It is an element of \(\mathcal{S}\) . For every \(\mathfrak{m}\) in \(\mathcal{S}\) , the field L/ \(\mathfrak{m}\) is a finite extension of K. We write \(\pi_{\mathfrak{m}}: \mathrm{L} \to \mathrm{L}/\mathfrak{m}\) for the projection and denote by \(\mathrm{G}_{\mathfrak{m}}\) the stabilizer of \(\mathfrak{m}\) in G, that is, the set of \(g \in \mathrm{G}\) such that \(g \cdot \mathfrak{m} = \mathfrak{m}\) . The K-algebra L/ \(\mathfrak{m}\) is endowed with an action of \(\mathrm{G}_{\mathfrak{m}}\) through the homomorphism \(\lambda_{\mathfrak{m}}\) from \(\mathrm{G}_{\mathfrak{m}}\) to \(\mathrm{Aut}_{\mathrm{K}}(\mathrm{L}/\mathfrak{m})\) that sends an element \(h\) of \(\mathrm{G}_{\mathfrak{m}}\) to the automorphism of L/ \(\mathfrak{m}\) deduced from \(\lambda(h)\) when passing to the quotients.

Let \(\mathcal{O}\) be the set of orbits of G in \(\mathcal{S}\) . Given an orbit \(\sigma \in \mathcal{O}\) , set \(\mathfrak{a}_{\sigma} = \bigcap_{\mathfrak{m} \in \sigma} \mathfrak{m}\) and \(\mathrm{L}_{\sigma} = \mathrm{L} / \mathfrak{a}_{\sigma}\) . Since \(\mathfrak{a}_{\sigma}\) is invariant under G, by passing to the quotients, the action of G on L defines a homomorphism \(\lambda_{\sigma}\) from G to \(\mathrm{Aut}_{\mathrm{K}}(\mathrm{L}_{\sigma})\) . Finally, denote by \(\pi_{\sigma}\) the canonical mapping from L to \(\mathrm{L}_{\sigma}\) .

Lemma 8. --- a) For every \(g \in G\) , \(\sigma \in \mathcal{O}\) , and \(\mathfrak{m} \in \sigma\) , we have

\[
[ \mathrm{L} _ {\sigma}: \mathrm{K} ] = \operatorname{Card} (\sigma) [ \mathrm{L} / \mathfrak {m}: \mathrm{K} ].
\]

\begin{enumerate}
\def\labelenumi{\alph{enumi})}
\setcounter{enumi}{1}
\item
  The mapping \(\pi : x \mapsto (\pi_{\sigma}(x))_{\sigma \in \mathcal{O}}\) is an isomorphism of (K,G)-algebras from L to \(\prod_{\sigma \in \mathcal{O}} \mathrm{L}_{\sigma}\) .
\item
  Denote by \(L^{G}\) (resp. \(L_{\sigma}^{G}\) ) the subalgebra of L (resp. \(L_{\sigma}\) ) of elements invariant under the action of G. Then \(\pi\) induces an isomorphism from \(L^{G}\) to \(\prod_{\sigma\in\mathcal{O}}L_{\sigma}^{G}\) .
\end{enumerate}

Since the algebra L is reduced and of finite degree, the intersection of the maximal ideals of L is reduced to 0 (VIII, p.~173, Corollary 2). Moreover, if m and \(m'\) are two distinct maximal ideals of L, then we have \(m + m' = L\) . By Proposition 10 of I, §8, No.~11, p.~110, the canonical mapping from L to \(\prod_{m \in S} L/m\) is an isomorphism, as is the canonical mapping from \(L/a_{\sigma}\) to \(\prod_{m \in \sigma} L/m\) for every \(\sigma \in O\) . Assertion a) follows. Since O is a partition of S, assertion b) follows; assertion c) is an immediate consequence of b).

Let us now fix an orbit \(\sigma \in \mathcal{O}\) and an element \(\mathfrak{m}\) of \(\sigma\) . Set \(F_{\mathfrak{m}} = \text{Coind}_{G_{\mathfrak{m}}}^{G}(L/\mathfrak{m})\) , and denote by \(\lambda_{F_{\mathfrak{m}}}\) the action of \(G\) on \(F_{\mathfrak{m}}\) . For any \(x \in L\) , denote by \(\overline{x}\) the mapping from \(G\) to \(L/\mathfrak{m}\) that sends \(g\) to \(\pi_{\mathfrak{m}}(gx)\) . By formula (15) of VIII, p.~305 and the definition of the action of \(G_{\mathfrak{m}}\) on \(L/\mathfrak{m}\) , it is immediate that \(\overline{x}\) belongs to \(F_{\mathfrak{m}}\) , and the mapping \(u: x \mapsto \overline{x}\) from \(L\) to \(F_{\mathfrak{m}}\) satisfies \(\lambda_{F_{\mathfrak{m}}}(g) \circ u = u \circ \lambda(g)\) for \(g \in G\) . In other words, \(u\) is a morphism of \((K, G)\) -algebras.

Lemma 9. --- The morphism \(u\) is surjective with kernel \(\mathfrak{a}_{\sigma}\) . The mapping deduced from \(u\) by passing to the quotient is an isomorphism of (K, G)-algebras from \(L_{\sigma}\) to \(F_{\mathfrak{m}}\) .

Since the kernel of \(\pi_{\mathfrak{m}}\) is equal to \(\mathfrak{m}\) , that of \(u\) is equal to \(\bigcap_{g\in G}g^{-1}.\mathfrak{m} = \mathfrak{a}_{\sigma}\) . To prove that \(u\) is surjective, it suffices to prove that the vector spaces \(\mathrm{L}_{\sigma} = \mathrm{L} / \mathfrak{a}_{\sigma}\) and \(\mathrm{F}_{\mathfrak{m}}\) have the same dimension over \(\mathrm{K}\) . Now, all ideals \(g\cdot \mathfrak{m}\) have the same codimension in \(\mathrm{L}\) , and, by Lemma 8, a),

\[
[ \mathrm{L} _ {\sigma}: \mathrm{K} ] = \operatorname{Card} (\sigma). [ \mathrm{L} / \mathfrak {m}: \mathrm{K} ].
\]

Moreover, by Lemma 7, b), we have

\[
[ \mathrm{F} _ {\mathfrak {m}}: \mathrm{K} ] = (\mathrm{G}: \mathrm{G} _ {\mathfrak {m}}) [ \mathrm{L} / \mathfrak {m}: \mathrm{K} ].
\]

Since we have \(\operatorname{Card}(\sigma)=(\mathrm{G}:\mathrm{G}_{\mathfrak{m}})\) , we have proved the equality \([L_{\sigma}:K]=[F_{m}:K]\) .

We now suppose, furthermore, that the homomorphism \(\lambda\) from G to \(\mathrm{Aut}_{\mathrm{K}}(\mathrm{L})\) is injective. We identify K with its image in L through the mapping \(\xi \mapsto \xi \cdot 1\) . Let \(\Omega\) be an algebraically closed extension of K. The set \(\mathcal{F}(\mathrm{G},\Omega)\) of mappings from G to \(\Omega\) coincides with the coinduced \((\Omega ,\mathrm{G})\) -algebra \(\mathrm{Coind}_{\{e\}}^{\mathrm{G}}(\Omega)\) . It is a free \(\Omega [\mathrm{G}]\) -module of rank 1. Let \(\mathcal{H}\) be the set of K-algebra homomorphisms from L to \(\Omega\) . We define a right action of G on \(\mathcal{H}\) by \((g,\chi)\mapsto \chi \circ \lambda (g)\) . The \(\Omega\) -algebra \(\mathcal{F}(\mathcal{H},\Omega)\) is then endowed with the structure of an \((\Omega ,\mathrm{G})\) -algebra deduced from the right action of G on \(\mathcal{H}\) .

Lemma 10. --- Let L be a (K, G)-algebra that is étale over K. The mapping \(\psi\) from \(\mathrm{L}_{(\Omega)}\) to \(\mathcal{F}(\mathcal{H},\Omega)\) characterized by the relation

\[
\psi (\xi \otimes x) = (\xi \chi (x)) _ {\chi \in \mathscr {H}}
\]

is an isomorphism of (K, G)-algebras.

Since L is étale, the mapping \(\psi\) is an isomorphism of \(\Omega\) -algebras (V, §6, No.~3, p.~30, Proposition 2 and V, §6, No.~3, p.~29, Proposition 1, c)). We have the relations

\[
\psi ((\operatorname{Id} \otimes \lambda (g)) (\xi \otimes x)) = (\xi (\chi \circ \lambda (g)) (x)) _ {\chi \in \mathscr {H}}
\]

for \(\xi \in \Omega\) , \(x \in \mathrm{L}\) , and \(g \in \mathrm{G}\) . So \(\psi\) is a morphism of \((\Omega, \mathrm{G})\) -algebras.

THEOREM 2. --- Let G be a finite group, and let L be a commutative K-algebra of finite degree endowed with an action of G given by an injective homomorphism \(\lambda\) from G to \(\operatorname{Aut}_{\mathrm{K}}(\mathrm{L})\) . The following properties are then equivalent:

\begin{enumerate}
\def\labelenumi{(\roman{enumi})}
\item
  There exist a subgroup H of G, a Galois extension E of K of finite degree, an isomorphism from H to \(\operatorname{Gal}(\operatorname{E}/\operatorname{K})\) , and an isomorphism of \((\operatorname{K},\operatorname{G})\) -algebras from L to \(\operatorname{Coind}_{\operatorname{H}}^{\operatorname{G}}(\operatorname{E})\) .
\item
  The algebra L is étale, and H is a homogeneous principal G-set (I, §5, No.~6, p.~60, Definition 7).
\item
  There exists an isomorphism of \((\Omega,\mathrm{G})\) -algebras \(\psi:\mathrm{L}_{(\Omega)}\to\mathcal{F}(\mathrm{G},\Omega)\) ; in other words, for every \(g\in G\) , the automorphism \(\psi\circ\lambda(g)_{(\Omega)}\circ\psi^{-1}\) of \(\Omega^{G}\) is equal to the automorphism
\end{enumerate}

\[
(x _ {h}) _ {h \in \mathrm{G}} \longmapsto (x _ {h g}) _ {h \in \mathrm{G}}.
\]

\begin{enumerate}
\def\labelenumi{(\roman{enumi})}
\setcounter{enumi}{3}
\item
  The algebra \(\mathrm{L}\) is reduced, and \(\mathrm{L}\) is a free \(\mathrm{K}[\mathrm{G}]\) -module of rank 1.
\item
  The algebra L is reduced, we have \(\operatorname{Card}(G) = [L : K]\) , and K is the subring of L of elements invariant under the action of G.
\item
  The algebra L is reduced, the group G acts transitively on the set of maximal ideals of L, and, for every maximal ideal m of L, the stabilizer \(G_{m}\) of m in G acts faithfully on L/m and admits K as subfield of invariants.
\end{enumerate}

(i)⇒(ii): Let E be a Galois extension of K of finite degree and \(\tau\) an isomorphism from H to \(\operatorname{Aut}_{\mathrm{K}}(\mathrm{E})\) . Let S be a system of representatives of the right cosets of G modulo H. The K-algebra \(F = \operatorname{Coind}_{\mathrm{H}}^{\mathrm{G}}(\mathrm{E})\) is isomorphic to \(\mathcal{F}(\mathrm{S}, \mathrm{E})\) (Lemma 7, a)); it is therefore étale. We denote by \(\lambda_{F}\) the action of G on F. Furthermore, let \(\psi\) be a K-algebra homomorphism from E to \(\Omega\) , and let \(\chi_{0}\) be the homomorphism \(f \mapsto \psi(f(e))\) from F to \(\Omega\) . Let \(g \in G\) be such that we have \(\chi_{0} \circ \lambda_{\mathrm{F}}(g) = \chi_{0}\) ; since \(\psi\) is injective, we then have \(f(g) = f(e)\) for every \(f \in F\) . In view of Lemma 7, a), we have \(g \in H\) , and by the formula

\[
f (h) = \tau (h) \cdot f (e),\tag{18}
\]

which holds for every \(h \in H\) , this can only happen if g = e. On the other hand, by Lemma 7, b) and Theorem 3 of V, §10, No.~6, p.~66, we have

\[
[ \mathrm{F}: \mathrm{K} ] = (\mathrm{G}: \mathrm{H}) [ \mathrm{E}: \mathrm{K} ] = (\mathrm{G}: \mathrm{H}) \operatorname{Card} (\mathrm{H}) = \operatorname{Card} (\mathrm{G}).
\]

The set \(\mathcal{H}\) of K-homomorphisms from F to \(\Omega\) has cardinal {[}F : K{]} because F is étale (V, §6, No.~5, p.~32, Proposition 4), so \(\mathrm{Card}(\mathcal{H}) = \mathrm{Card}(\mathrm{G})\) . Since the stabilizer of \(\chi_0\) in G is equal to \(\{e\}\) by the above, \(\mathcal{H}\) is a homogenous principal G-set.

(ii)⇒(iii): Suppose that L is étale and that H is a homogenous principal G-set. By Lemma 10, the \((\Omega,\mathrm{G})\) -algebras \(\mathrm{L}_{(\Omega)}\) and \(\mathcal{F}(\mathcal{H},\Omega)\) are isomorphic. Since H is a homogenous principal G-set, the \((\Omega,\mathrm{G})\) -algebras \(\mathcal{F}(\mathcal{H},\Omega)\) and \(\mathcal{F}(\mathrm{G},\Omega)\) are isomorphic.

(iii)⇒(iv): Suppose that property (iii) holds. Then \(L_{(\Omega)}\) is a free module of rank 1 over the algebra \(\Omega[G]\) ; the latter can be canonically identified with \(K[G]_{(\Omega)}\) . We then apply Theorem 3 of VIII, p.~37.

The implication \((\mathrm{iv})\Rightarrow (\mathrm{v})\) is immediate.

Let us prove the implication \((v)\Rightarrow(vi)\) . The algebra L is reduced. By Lemma 8, c), the group G acts transitively on the set S of maximal ideals of L. Let m be an element of S. By Lemma 9, since \(\bigcap_{n\in S}n=\{0\}\) , the algebra L is isomorphic to the algebra \(\operatorname{Coind}_{G_{m}}^{G}(L/m)\) . The algebra of invariants of \(G_{m}\) in L/m therefore coincides with K by Lemma 7, c). Hence the homomorphism \(\lambda_{m}\) from \(G_{m}\) to \(\operatorname{Gal}((L/m)/K)\) is surjective. By Lemma 7, we moreover have

\[
\operatorname{Card} (\mathrm{G}) = [ \mathrm{L}: \mathrm{K} ] = (\mathrm{G}: \mathrm{G} _ {\mathfrak {m}}) [ \mathrm{L} / \mathfrak {m}: \mathrm{K} ].
\]

So \(\operatorname{Card}(\mathrm{G}_{\mathfrak{m}}) = [\mathrm{L} / \mathfrak{m}:\mathrm{K}]\) , and the homomorphism \(\lambda_{\mathfrak{m}}\) is injective.

It remains to prove the implication \((\mathrm{vi})\Rightarrow(\mathrm{i})\) . Let m be a maximal ideal of L. By Lemma 9, the algebra L is isomorphic to the algebra \(\operatorname{Coind}_{G_{m}}^{G}(L/m)\) as a \((K,G)\) -algebra. Since \(G_{m}\) acts faithfully in L/m and admits K as subfield of invariants, the group homomorphism \(\lambda_{m}\) defines an isomorphism from \(G_{m}\) to \(\operatorname{Gal}((L/m)/K)\) .

Remarks. --- 2) In Theorem 2, we can replace the assumption that \(\Omega\) is algebraically closed with the assumption that \(\Omega\) is separably closed. Indeed, if \(L\) is étale, then the image of every K-algebra homomorphism from \(L\) to \(\Omega\) is a separable extension of \(K\) .

\begin{enumerate}
\def\labelenumi{\arabic{enumi})}
\setcounter{enumi}{2}
\tightlist
\item
  The normal basis theorem (V, §10, No.~9, p.~73, Theorem 6) is a specific case of the implication \((\mathrm{i})\Rightarrow (\mathrm{iv})\) in Theorem 2.
\end{enumerate}

DEFINITION 1. --- Let G be a finite group, and let L be a nonzero commutative algebra of finite degree over K endowed with the structure of a (K,G)-algebra, where the action of G is given by an injective homomorphism from G to the group \(\operatorname{Aut}_{\mathrm{K}}(\mathrm{L})\) . We say that L is a Galois algebra with group G if it has the equivalent properties (i) through (vi) of Theorem 2.

Remarks. --- 4) Suppose that L is an extension of K endowed with an action \(\lambda\) of G. Then L is a Galois algebra over K if and only if L is a Galois extension of K and \(\lambda\) is an isomorphism from G to the Galois group of L over K.

\begin{enumerate}
\def\labelenumi{\arabic{enumi})}
\setcounter{enumi}{4}
\tightlist
\item
  Suppose that the group G is abelian. If G acts faithfully and transitively on a set X, then the stabilizer of any point of X is reduced to the identity (because the stabilizers of the points of X are all equal and their intersection is reduced to the identity element of G). Consequently, properties (i) through (vi) of Theorem 2 are also equivalent to the following:
\end{enumerate}

\begin{enumerate}
\def\labelenumi{(\roman{enumi})}
\setcounter{enumi}{6}
\tightlist
\item
  The algebra \(\mathrm{L}\) is étale, and \(\mathrm{G}\) acts faithfully and transitively on \(\mathcal{H}\) .
\end{enumerate}

\begin{enumerate}
\def\labelenumi{\arabic{enumi})}
\setcounter{enumi}{5}
\tightlist
\item
  By V, §10, No.~1, p.~75, Lemma 5; V, §10, No.~1, p.~76, Proposition 12; and VIII, p.~137, Proposition 3, a Galois algebra is commutative and semisimple.
\end{enumerate}

Examples. --- 1) Let n be a strictly positive integer, prime to the characteristic exponent of K. Suppose that the group \(\mu_{n}\) of n-th roots of unity in K has order n.~Then for every divisor d of n, the group \(\mu_{d}\) of d-th roots of unity has order d.~Let a be a nonzero element of K, let L be the algebra \(\mathrm{K}[\mathrm{X}]/(X^{n}-a)\) , and let x be the class of X in L. The sequence \((1,x,\ldots,x^{n-1})\) is a basis of the vector space L over the field K, and we have \(x^{n}=a\) . Moreover, the polynomial

\(X^{n}-a\) is prime to its derivative \(nX^{n-1}\) , so the algebra L is étale (V, §7, No.~2, p.~37, Proposition 3). For every \(\zeta\) in \(\mu_{n}\) , the automorphism \(P(X)\mapsto P(\zeta X)\) of the ring K{[}X{]} defines, by passing to the quotients, an endomorphism \(\lambda(\zeta)\) of L because we have \((\zeta X)^{n}-a=X^{n}-a\) ; it is an automorphism. We have

\[
\lambda (\zeta) x ^ {i} = \zeta^ {i} x ^ {i}\tag{19}
\]

for \(0 \leqslant i < n\) . The mapping \(\lambda : \zeta \mapsto \lambda(\zeta)\) is an injective homomorphism from \(\mu_{n}\) to \(\operatorname{Aut}_{\mathrm{K}}(\mathrm{L})\) , and the ring of invariants of the group \(\lambda(\mu_{n})\) in L is equal to K·1. Since the cardinal of \(\mu_{n}\) is equal to \(n = [L : K]\) , the algebra L endowed with the action of \(\lambda\) is a Galois algebra (VIII, p.~308, Theorem 2, (v)).

Let \(r\) be the least strictly positive integer such that \(a^r\) belongs to \(\mathrm{K}^{*n}\) ; it divides \(n\) , and there exists an element \(b\) of \(\mathrm{K}^*\) such that \(a = b^{n / r}\) . Then (V, §11, No.~8, p.~91, Remark) the polynomial \(\mathrm{X}^r - b\) is irreducible, and we have \(\mathrm{X}^n - a = \prod_{\zeta \in \mu_{n / r}} (\mathrm{X}^r - \zeta b)\) . Let \(\mathrm{E}\) be the field \(\mathrm{K}[\mathrm{Y}] / (\mathrm{Y}^r - b)\) , and let \(y\) be the class of \(Y\) in \(\mathrm{E}\) . There exists an isomorphism \(\theta\) from \(\mu_{n / r}\) to \(\operatorname{Gal}(\mathrm{E} / \mathrm{K})\) , characterized by the relation \(\theta(\xi)(y) = \xi y\) (V, §11, No.~8, p.~91, Example 3). We then verify that the Galois algebra \(L\) is isomorphic to the \((K, \mu_n)\) -algebra \(\operatorname{Coind}_{\mu_{n / r}}^{\mu_n}(E)\) .

\begin{enumerate}
\def\labelenumi{\arabic{enumi})}
\setcounter{enumi}{1}
\tightlist
\item
  Now suppose that the field K has characteristic \(p \neq 0\) . Let c be an element of K. The polynomial \(f = X^{p} - X - c\) is prime to its derivative \(f' = -1\) , so the algebra \(L = K[X]/(f)\) is étale (V, §7, No.~2, p.~37, Proposition 3). We denote the image of X in L by x; we have the relation \(x^{p} = x + c\) . The sequence \((1, x, \ldots, x^{p-1})\) is a basis of L viewed as a vector space over K.
\end{enumerate}

Let P be the additive group of the prime subfield of K; it is a cyclic group of order p, generated by the unit element 1 of K. For every j in P, we have \(j^{p} = j\) (V, §1, No.~3, p.~4, formula (4)) and therefore \(f(\mathrm{X} + j) = f(\mathrm{X})\) . Hence there exists an automorphism \(\gamma(j)\) of the algebra L characterized by the relation \(\gamma(j)(x) = x + j\) ; moreover, the resulting mapping \(\gamma\) is an injective homomorphism from P to \(\operatorname{Aut}_{\mathrm{K}}(\mathrm{L})\) .

Let \(\Omega\) be an algebraically closed extension of \(K\) , and let \(\xi\) be a root of the polynomial \(f\) in \(\Omega\) . We have \(\xi^p = \xi + c\) , hence

\[
\mathrm{X} ^ {p} - \mathrm{X} - c = (\mathrm{X} ^ {p} - \xi^ {p}) - (\mathrm{X} - \xi) = (\mathrm{X} - \xi) ^ {p} - (\mathrm{X} - \xi) = \prod_ {j \in \mathrm{P}} (\mathrm{X} - \xi - j)
\]

by V, §12, No.~1, p.~94, formula (1). For every \(j\) in P, there exists a unique algebra homomorphism \(\chi_j: \mathrm{L} \to \Omega\) that sends \(x\) to \(\xi + j\) ; moreover, every homomorphism from L to \(\Omega\) is one of the \(\chi_j\) , and we have the relation \(\chi_{j}=\chi_{0}\circ\gamma(j)\) . The algebra L endowed with \(\gamma\) has property (ii) of Theorem 2 of VIII, p.~308 and is therefore a Galois algebra over K.

To describe the structure of L, we must distinguish between two cases:

\begin{enumerate}
\def\labelenumi{\alph{enumi})}
\item
  We have \(\xi\notin K\) . Then the polynomial \(f(X)\) is irreducible in \(K[X]\) (V, §11, No.~9, p.~93, Example 3). In this case, L is a cyclic extension of K of degree p, and \(\gamma\) is an isomorphism from P to \(\operatorname{Gal}(L/K)\) .
\item
  We have \(\xi \in \mathrm{K}\) . Then the mapping \(\psi : y \mapsto (\chi_j(y))_{j \in \mathrm{P}}\) is an isomorphism from the algebra \(\mathrm{L}\) to the product algebra \(\mathrm{K}^{\mathrm{P}}\) ; moreover, \(\psi \circ \gamma(k) \circ \psi^{-1}\) is the automorphism \((x_j)_{j \in \mathrm{P}} \mapsto (x_{j+k})_{j \in \mathrm{P}}\) of \(\mathrm{K}^{\mathrm{P}}\) for every \(k \in \mathrm{P}\) .
\end{enumerate}

